\section{Introduction}
\subsection{Purpose} % ------------------------------------------------------------------
The \textit{Best Bike Paths} (\textit{BBP}) system has one objective in mind: support the ever-growing cycling community by providing a platform for recording, browsing, and sharing optimal bike routes. By integrating automated data acquisition from mobile sensors with manual user contributions, BBP builds a reliable inventory of safe and smooth paths. It enhances the biking experience by offering detailed performance statistics and real-time meteorological information on the go. Furthermore, BBP promotes safety by automatically identifying physical road hazards through community-verified data, which non-registered users can also benefit from. By prioritizing routes where speed limits and traffic density are bike-friendly, it promotes a safer environment for riders of all skill levels. Through its dual-mode insertion and analytical scoring, BBP becomes a comprehensive tool for both personal fitness and community navigation.

\subsubsection{Goals} % ------------------------------------------------------------------
\begin{itemize}
    \item G1: Allow registered users to record, store and track their cycling activities in order to have a history of their activities. (!!)
    \item G2: Allow registered users to record and store their cycling activities in order to enrich and update the paths inventory. (!!)
    \item G3: Allow users to continuously update information about bike paths.
    \item G4: Display statistics about the trip, including the total distance, average speed and performance metrics.
    \item G5: Enrich trip data with meteorological information retrieved from external services.
    \item G6: Support registered users to manually insert information about a bike path, including the name of the street, its status and the presence of obstacles along the way, via a virtual map. (!!) 
    \item G7: Make data available to all users, registered or not.
\end{itemize}

\subsection{Scope} % ------------------------------------------------------------------

\subsubsection{World Phenomena}
\begin{itemize}
    \item WP1: Users physically drive their bike from one location to another over a period of time.
    \item WP2: Mobile devices record changes of their sensors, like GPS coordinates \footnote{GPS coordinates are registered by the device only if the user enables it. We assume that the device location is always on.}, orientation, speed and acceleration, independently from the system. 
    \item WP3: Weather conditions, temperature and wind intensity change in a specific location over a period of time.
    \item WP4: Road conditions change in a specific location over a period of time.
    \item WP5: Users encounter obstacles along their way.
    \item WP6: Users ride their bikes in trafficked roads.
    \item WP7: Users develop opinions over certain paths.
    \item WP8: Users decide to accept the automatic tracking of the system.
\end{itemize}

\subsubsection{Shared Phenomena}

\paragraph{World Controlled Shared Phenomena}
\begin{itemize}
    \item SP1: Users register to the service via the specific form.
    \item SP2: Users log in via the specific form.
    \item SP3: Users start the recording of their biking session.
    \item SP4: External services send weather data to the system, if available, under request.
    \item SP5: Registered users manually insert bike path status information via the UI.
    \item SP6: Registered users correct automatically acquired data about the bike path status via the UI.
    \item SP7: Registered users confirm automatically acquired data about the bike path status via the UI.
    \item SP8: Users publish acquired information about paths.
    \item SP9: Users browse possible paths.
    \item SP10: Users plan their biking trip by entering the departure and destination point via the UI.
    \item SP11: Users accept the automatic tracking of the system.
\end{itemize}

\paragraph{Machine Controlled Shared Phenomena}
\begin{itemize}
    \item SP12 The system records and stores cycling activities, including metrics, such as total distance covered and average speed. 
    \item SP13 Mobile devices send GPS coordinates and sensor data from the gyroscope and accelerometer to the system, if automatic tracking accepted.
    \item SP14 The system sends a request to external services to query weather information.
    \item SP15 The system collects biking path information.
    \item SP16 The system computes a path score based on the available information relative to the path.
    \item SP17 The system visualizes the available bike paths connecting the departure point to the destination point based on their rank.
    \item SP18 The system signals to the user obstacles along paths.
\end{itemize}

\subsection{Definitions, Acronyms, Abbreviations} % ------------------------------------------------------------------
\subsubsection{Definitions}


\subsubsection{Acronyms}
\begin{itemize}
    \item BBP: Best Bike Paths.
    \item GPS: Global Positioning System.
    \item UI: User Interface.
\end{itemize}

\subsubsection{Abbreviations}
\begin{itemize}
    \item G: goal.
    \item WP: World Phenomena.
    \item SP: Shared Phenomena.
\end{itemize}

\subsection{Revision history}
\begin{itemize}
    \item Version 1.0: 23/12/2025
\end{itemize}

\subsection{Reference Documents} % ------------------------------------------------------------------
The assignment for this document and all the information included refer to the following
documentation:
\begin{itemize}
    \item The specification for the 2025/26 Requirement Engineering and Design Project for the Software Engineering II course.
    \item The slides on the WeBeep page of the Software Engineering II course.
\end{itemize}

\subsection{Document Structure} % ------------------------------------------------------------------
\begin{enumerate}
    \item Introduction: a brief description of the project that follows from the analysis of the specification provided, it highlights the goals of the product and its scope.
    \item Overall Description: a more detailed description of the product that comprehends also real world scenarios and analyses the system from a technical point of view.
    \item Specific Requirements: a detailed list of all the requirements needed in order to satisfy the goals, and their relationship with domain assumptions.
    \item Formal analysis using Alloy: a formal description of a portion of real world phenomena analysed through Alloy in order to highlight inconsistencies between assumptions and phenomena previously modelled.
    \item Effort spent: the time spent by each member of the group to contribute to the realization of this document.
    \item References: reference to documents or tools used for writing this document.
\end{enumerate}
